\documentclass[11pt,fontset=fandol]{ctexart}
\usepackage[a4paper,margin=23mm]{geometry}
\usepackage{amsmath,amssymb,mathtools,bm,mathrsfs,esint}
\usepackage{graphicx,longtable,booktabs,array,calc}
\usepackage[normalem]{ulem}
\usepackage{xcolor,hyperref}
\hypersetup{hidelinks}
\setcounter{secnumdepth}{-1}
\newcounter{none}
\providecommand{\tightlist}{\setlength{\itemsep}{0pt}\setlength{\parskip}{0pt}}
\providecommand{\passthrough}[1]{#1}
\usepackage{amssymb,mathrsfs,esint,bm,newunicodechar}
\allowdisplaybreaks
\setlength{\emergencystretch}{3em}
\providecommand{\llangle}{\langle\!\langle}
\providecommand{\rrangle}{\rangle\!\rangle}
\newunicodechar{∞}{\ensuremath{\infty}}
\newunicodechar{∈}{\ensuremath{\in}}
\newunicodechar{∑}{\ensuremath{\sum}}
\newunicodechar{∫}{\ensuremath{\int}}
\newunicodechar{≤}{\ensuremath{\leq}}
\newunicodechar{≥}{\ensuremath{\geq}}
\newunicodechar{≠}{\ensuremath{\neq}}
\newunicodechar{→}{\ensuremath{\to}}
\newunicodechar{⇒}{\ensuremath{\Rightarrow}}
\newunicodechar{α}{\ensuremath{\alpha}}
\newunicodechar{β}{\ensuremath{\beta}}
\newunicodechar{γ}{\ensuremath{\gamma}}
\newunicodechar{λ}{\ensuremath{\lambda}}
\newunicodechar{φ}{\ensuremath{\varphi}}
\newunicodechar{ε}{\ensuremath{\varepsilon}}
\newunicodechar{σ}{\ensuremath{\sigma}}
\newunicodechar{θ}{\ensuremath{\theta}}
\newunicodechar{δ}{\ensuremath{\delta}}
\newunicodechar{η}{\ensuremath{\eta}}
\newunicodechar{π}{\ensuremath{\pi}}
\newunicodechar{ξ}{\ensuremath{\xi}}
\newunicodechar{Σ}{\ensuremath{\Sigma}}
\newunicodechar{Δ}{\ensuremath{\Delta}}
\newunicodechar{∂}{\ensuremath{\partial}}
\newunicodechar{∀}{\ensuremath{\forall}}
\newunicodechar{∃}{\ensuremath{\exists}}
\newunicodechar{√}{\ensuremath{\surd}}
\newunicodechar{−}{\ensuremath{-}}
\newunicodechar{μ}{\ensuremath{\mu}}
\newunicodechar{ρ}{\ensuremath{\rho}}
\newunicodechar{τ}{\ensuremath{\tau}}
\newunicodechar{ω}{\ensuremath{\omega}}
\newunicodechar{Ω}{\ensuremath{\Omega}}
\newunicodechar{χ}{\ensuremath{\chi}}
\newunicodechar{ψ}{\ensuremath{\psi}}
\newunicodechar{Γ}{\ensuremath{\Gamma}}
\newunicodechar{Φ}{\ensuremath{\Phi}}
\newunicodechar{Λ}{\ensuremath{\Lambda}}
\newunicodechar{ℝ}{\ensuremath{\mathbb{R}}}
\newunicodechar{ℤ}{\ensuremath{\mathbb{Z}}}
\newunicodechar{ℕ}{\ensuremath{\mathbb{N}}}
\newunicodechar{△}{\ensuremath{\triangle}}
\newunicodechar{⊗}{\ensuremath{\otimes}}
\begin{document}
\section{2013年杭州电子科技大学601数学分析考研真题}\label{ux5e74ux676dux5ddeux7535ux5b50ux79d1ux6280ux5927ux5b66601ux6570ux5b66ux5206ux6790ux8003ux7814ux771fux9898}

来源：\texttt{HDU/601数学分析/2013/2013年杭州电子科技大学601数学分析考研真题.pdf}。原
PDF 页序：1, 2, 3。

\section{杭州电子科技大学}\label{ux676dux5ddeux7535ux5b50ux79d1ux6280ux5927ux5b66}

\section{2013年攻读硕士学位研究生入学考试}\label{ux5e74ux653bux8bfbux7855ux58ebux5b66ux4f4dux7814ux7a76ux751fux5165ux5b66ux8003ux8bd5}

\section{《数学分析》试题}\label{ux6570ux5b66ux5206ux6790ux8bd5ux9898}

\begin{quote}
助手整理：按同题参考答案中可相互核对的表达式恢复 OCR
损坏的公式。二（3）的目标仍为原题的 \(xu_{xx}+yu_{xy}\)。
\end{quote}

（试题共四大题，共3页，总分150分）

姓名 报考专业 准考证号

【所有答案必须写在答题纸上，做在试卷或草稿纸上无效！】

\subsection{一、填空题（本大题共8小题，每小题5分，本大题共40分）}\label{ux4e00ux586bux7a7aux9898ux672cux5927ux9898ux51718ux5c0fux9898ux6bcfux5c0fux98985ux5206ux672cux5927ux9898ux517140ux5206}

1.求
\(\lim _ { x \to 1 } \left( { \frac { 2 x } { 1 + x } } \right) ^ { \frac { 2 x } { 1 - x } }\)

\begin{enumerate}
\def\labelenumi{\arabic{enumi}.}
\setcounter{enumi}{1}
\item
  \(\int \sin { \frac { 1 } { x } } \cdot { \frac { d x } { x ^ { 2 } } } =\)
\item
  设 \(z=e^{\sin(xy)}\)，则 \(dz=\)\_\_\_\_。
\item
  假设
  \(\int_0^1dx\int_{x^2}^xf(x,y)\,dy=\int_0^1dy\int_y^{\varphi(y)}f(x,y)\,dx\)，则
  \(\varphi(y)=\)\_\_\_\_。
\end{enumerate}

5.设 \(y = ( 1 + \sin x ) ^ { x }\) ，则
\(\left. d y \right| _ { x = \pi } =\)

\begin{enumerate}
\def\labelenumi{\arabic{enumi}.}
\setcounter{enumi}{5}
\item
  设函数
  \(f(x)=\begin{cases}\sin(ax)/x,&x<0,\\2-e^x,&x\geq0,\end{cases}\)，若
  \(f\) 为连续函数，则 \(a=\)\_\_\_\_。
\item
  已知 \(f(x)\) 和 \(\varphi(x)\) 是 \((-\infty,+\infty)\)
  上的连续函数，且 \(\int_1^{x^2}f(t)\,dt=\int_1^x\varphi(t)\,dt\)
  成立，则 \(\varphi(x)=\)\_\_\_\_。
\end{enumerate}

8.计算无穷级数的和
\(\sum _ { n = 1 } ^ { \infty } n { \Bigg ( } { \frac { 1 } { 2 } } { \Bigg ) } ^ { n }\)

\subsection{二、计算题（本大题共4小题，本大题共30分）}\label{ux4e8cux8ba1ux7b97ux9898ux672cux5927ux9898ux51714ux5c0fux9898ux672cux5927ux9898ux517130ux5206}

1.求极限
\(\lim _ { x \to 0 } \left( { \frac { 1 } { \ln ( 1 + x ) } } - { \frac { 1 } { x } } \right)\)
（本小题5分）

2．计算定积分
\(\int _ { 0 } ^ { \ln 2 } { \sqrt { e ^ { x } - 1 } } d x\)
．（本小题5分）

3.设
\(u = y f { \left( { \frac { x } { y } } \right) } + x g { \left( { \frac { y } { x } } \right) } ,\)
f与g具有二阶连续导数，计算\(x \frac { \partial ^ { 2 } u } { { \partial x } ^ { 2 } } + y \frac { \partial ^ { 2 } u } { { \partial x } { \partial y } }\)
（本小题10分）

4.求圆锥面 \(z ^ { 2 } = a ^ { 2 } ( x ^ { 2 } + y ^ { 2 } )\) 截圆柱面
\(x ^ { 2 } + y ^ { 2 } = 2 y\) 所得有界部分立体的体积．（本小题10分）

\subsection{三、综合题（本大题共3小题，每小题10分，本大题共30分）}\label{ux4e09ux7efcux5408ux9898ux672cux5927ux9898ux51713ux5c0fux9898ux6bcfux5c0fux989810ux5206ux672cux5927ux9898ux517130ux5206}

1．讨论函数 \(f ( x ) = \ln ( 1 + x ^ { 2 } )\)
的增减性、凹凸性、并求极值与拐点.

2.讨论级数
\(\sum _ { n = 1 } ^ { \infty } { \frac { 1 } { x ^ { n } } } \sin { \frac { \pi } { 2 ^ { n } } }\)
的敛散性，并求出收敛域.

\begin{enumerate}
\def\labelenumi{\arabic{enumi}.}
\setcounter{enumi}{2}
\tightlist
\item
  设 \(C^+\) 为从 \(A(1,0)\) 到 \(B(-1,0)\) 的上半圆
  \(y=\sqrt{1-x^2}\)，\(-1\leq x\leq1\)。计算
  \(I=\int_{C^+}(-2xe^{-x^2}\sin y)\,dx+(e^{-x^2}\cos y+x^4)\,dy\)。
\end{enumerate}

\subsection{四、证明题（本大题共5小题，每小题10分，本大题共50分）}\label{ux56dbux8bc1ux660eux9898ux672cux5927ux9898ux51715ux5c0fux9898ux6bcfux5c0fux989810ux5206ux672cux5927ux9898ux517150ux5206}

1.已知x\textgreater0，证明：
\(x - \frac { x ^ { 2 } } { 2 } < \ln ( 1 + x ) < x\)

2.设函数列 \(\left\{ f _ { n } ( x ) \right\}\)
在区间1上一致收敛于f(x)，且 \(f _ { n } ( x )\) 在区间1上一致连续
\(( n \in N )\) .证明：f(x)在I上也一致连续.

3．设函数

\[
f(x,y)=\begin{cases}\dfrac{x^2y^2}{(x^2+y^2)^{3/2}},&x^2+y^2\neq0,\\0,&x^2+y^2=0.\end{cases}
\]

证明，f(x,y)在(0,0)处连续但不可微.

4.证明，当p\textgreater0时，积分
\(\int _ { 0 } ^ { + \infty } x ^ { p - 1 } e ^ { - x }\) dx收敛.

\begin{enumerate}
\def\labelenumi{\arabic{enumi}.}
\setcounter{enumi}{4}
\tightlist
\item
  设函数 \(f\) 在 \([0,1]\) 上可导，且
  \(\int_0^1xf(x)\,dx=f(1)\)。证明存在 \(\xi\in(0,1)\)，使
  \(f'(\xi)=-f(\xi)/\xi\)。
\end{enumerate}
\end{document}
